\newpage
\fancyhead[L]{\songti\zihao{-5}\cquHeader}
\fancyhead[R]{\songti\zihao{-5}附录A：Mordell 主定理的详细推导}
\addcontentsline{toc}{section}{附录A：Mordell 主定理的详细推导}
\section*{附录A：Mordell 主定理的详细推导}
\zihao{5}

\subsection*{A.1 在 \(\mathbb{Z}[\sqrt d]\) 中分解}

由
\[
y^2=x^3+d
\]
得到
\[
y^2-d=x^3.
\]
在
\[
D=\mathbb{Z}[\sqrt d]
\]
中，左边分解为
\[
(y+\sqrt d)(y-\sqrt d)=x^3.
\]
转为主理想：
\[
\langle y+\sqrt d\rangle\langle y-\sqrt d\rangle=\langle x\rangle^3.
\]
记
\[
I=\langle y+\sqrt d\rangle,\qquad
J=\langle y-\sqrt d\rangle,\qquad
L=\langle x\rangle.
\]
则
\[
IJ=L^3.
\]

\subsection*{A.2 证明 \(I,J\) 互素}

需要证明
\[
\gcd(I,J)=1.
\]
手算证明思路如下：若某个素理想同时整除 \(I\) 和 \(J\)，则它同时整除
\[
(y+\sqrt d)+(y-\sqrt d)=2y
\]
和
\[
(y+\sqrt d)-(y-\sqrt d)=2\sqrt d.
\]
因此公共因子只能来自 \(2\)、\(y\)、\(d\) 的公共约束。

论文中用到两个辅助事实：第一，\(y\) 与 \(d\) 互素；第二，在定理假设下，任何解都满足 \(x\) 为奇数。

\(x\) 为奇数可用模 \(8\) 检查：若 \(x\) 偶，则
\[
x^3\equiv 0\pmod 8.
\]
而平方模 \(8\) 只可能是
\[
0,1,4.
\]
对 \(d\equiv 2,3\pmod 4\) 的相关剩余类，
\[
y^2-d\equiv x^3\pmod 8
\]
会矛盾。因此 \(x\) 必为奇数。

Lean 形式化中没有采用完全相同的纸笔证明，而是直接通过范数限制计算两个理想的 gcd。

\subsection*{A.3 理想下降}

已经有
\[
IJ=L^3,\qquad \gcd(I,J)=1.
\]
由 Dedekind domain 中的理想下降定理，存在理想 \(A,B\) 使得
\[
I=A^3,\qquad J=B^3.
\]
特别地，
\[
\langle y+\sqrt d\rangle=A^3.
\]
因为左边是主理想，所以 \(A^3\) 是主理想。

\subsection*{A.4 类数条件推出 \(A\) 主理想}

在类群中，
\[
[A]^3=1.
\]
假设
\[
3\nmid h(D),
\]
即
\[
\gcd(3,h(D))=1.
\]
由有限群理论可得
\[
[A]=1.
\]
所以 \(A\) 是主理想：
\[
A=\langle\alpha\rangle
\]
对某个
\[
\alpha\in D.
\]
于是
\[
\langle y+\sqrt d\rangle=\langle\alpha^3\rangle.
\]
两个主理想相等意味着生成元相差一个单位：
\[
y+\sqrt d=u\alpha^3,\qquad u\in D^\times.
\]

\subsection*{A.5 单位都是三次方}

对
\[
D=\mathbb{Z}[\sqrt d],\qquad d<0,
\]
单位可以通过范数方程刻画：
\[
N(a+b\sqrt d)=a^2-db^2=1.
\]
当 \(d<-1\) 时，单位只有
\[
\pm 1.
\]
当 \(d=-1\) 时，单位为
\[
\{\pm1,\pm i\}.
\]
这两个单位群的阶都是 \(2\) 的幂，因此三次幂映射
\[
u\mapsto u^3
\]
是单位群的双射。换句话说，每个单位都是三次方。

所以可以把 \(u\) 吸收到 \(\alpha\) 中，不妨设
\[
y+\sqrt d=\alpha^3.
\]

\subsection*{A.6 展开 \(\alpha^3\)}

写
\[
\alpha=a+b\sqrt d,\qquad a,b\in\mathbb{Z}.
\]
先平方：
\[
(a+b\sqrt d)^2=a^2+2ab\sqrt d+b^2d.
\]
再乘一次：
\[
(a+b\sqrt d)^3=(a^3+3ab^2d)+(3a^2b+b^3d)\sqrt d.
\]
也可写成
\[
(a+b\sqrt d)^3=a(a^2+3b^2d)+b(3a^2+b^2d)\sqrt d.
\]
由于
\[
y+\sqrt d=(a+b\sqrt d)^3.
\]
比较 \(\sqrt d\) 的系数，得到
\[
1=b(3a^2+b^2d).
\]
这是整数中的分解，所以
\[
b=\pm1.
\]
又因为 \(b^2=1\)，上式给出
\[
3a^2+d=b.
\]
因此
\[
d=b-3a^2.
\]
即
\[
d=1-3a^2
\qquad\text{或}\qquad
d=-1-3a^2.
\]
令
\[
m=a,
\]
得到
\[
d=-3m^2\pm 1.
\]
再比较有理部分，得到
\[
y=a(a^2+3b^2d).
\]
由于 \(b^2=1\)，所以
\[
y=a(a^2+3d)=m(m^2+3d)=m(3d+m^2).
\]
这就是主定理中的 \(y\) 公式。

\subsection*{A.7 推出 \(x\) 公式}

对
\[
y+\sqrt d=(a+b\sqrt d)^3
\]
取范数：
\[
N(y+\sqrt d)=N(a+b\sqrt d)^3.
\]
左边是
\[
N(y+\sqrt d)=y^2-d.
\]
由 Mordell 方程，
\[
y^2-d=x^3.
\]
右边是
\[
N(a+b\sqrt d)^3=(a^2-db^2)^3.
\]
因为 \(b^2=1\)，所以
\[
x^3=(a^2-d)^3.
\]
整数中立方映射单射，因此
\[
x=a^2-d=m^2-d.
\]

\subsection*{A.8 反方向验证}

反方向更简单。若
\[
d=1-3m^2
\qquad\text{或}\qquad
d=-1-3m^2,
\]
令
\[
x=m^2-d,\qquad y=m(3d+m^2),
\]
则直接代入可证
\[
y^2=x^3+d.
\]
统一地，令
\[
s\in\{1,-1\},\qquad d=s-3m^2.
\]
设
\[
t=m^2.
\]
则
\[
3d+m^2=3s-8t,
\]
\[
m^2-d=4t-s.
\]
于是
\[
y^2=m^2(3d+m^2)^2=t(3s-8t)^2.
\]
展开：
\[
t(3s-8t)^2=t(9s^2-48st+64t^2).
\]
因为 \(s^2=1\)，所以
\[
y^2=9t-48st^2+64t^3.
\]
另一方面，
\[
x^3+d=(4t-s)^3+s-3t.
\]
展开：
\[
(4t-s)^3+s-3t
=64t^3-48st^2+12s^2t-s^3+s-3t.
\]
由于 \(s^2=1\) 且 \(s^3=s\)，得到
\[
x^3+d=64t^3-48st^2+12t-s+s-3t.
\]
即
\[
x^3+d=64t^3-48st^2+9t.
\]
这与 \(y^2\) 相同，因此
\[
y^2=x^3+d.
\]
Lean 中这个验证由 \texttt{ring} tactic 直接完成。

\newpage
\fancyhead[L]{\songti\zihao{-5}\cquHeader}
\fancyhead[R]{\songti\zihao{-5}附录B：Lean 4 代码片段}
\addcontentsline{toc}{section}{附录B：Lean 4 代码片段}
\section*{附录B：Lean 4 代码片段}
\zihao{5}

本附录用于整理正文中涉及的 Lean 4 形式化代码。正式定稿时，可将已经通过编译的核心定义、引理和定理放置在本处，并注明其所在 Lean 项目文件路径。

\begin{verbatim}
-- 示例：Mordell 方程整数解谓词
def IsMordellSolution (d x y : Int) : Prop :=
  y ^ 2 = x ^ 3 + d
\end{verbatim}

\newpage
\fancyhead[L]{\songti\zihao{-5}\cquHeader}
\fancyhead[R]{\songti\zihao{-5}附录C：AI 实验记录表}
\addcontentsline{toc}{section}{附录C：AI 实验记录表}
\section*{附录C：AI 实验记录表}
\zihao{5}

\begin{table}[htbp]
\centering
\caption*{AI 辅助 Lean 形式化实验记录表模板}
\small
\begin{tabular}{ccccc}
\toprule
任务编号 & 任务类型 & 模型 & 是否编译通过 & 主要问题\\
\midrule
T1 & 代数恒等式 & 待填写 & 待填写 & 待填写\\
T2 & 命题建模 & 待填写 & 待填写 & 待填写\\
T3 & 报错修复 & 待填写 & 待填写 & 待填写\\
\bottomrule
\end{tabular}
\end{table}
